\section{Related Work}

\paragraph{Duplex Models.}
Recent work in spoken dialogue systems (SDMs) increasingly draws on human conversational behaviors. Building on insights from human conversation, recent SDMs have progressed toward duplex capabilities—systems that listen and speak concurrently. SDMs are commonly built in two modes. Half-duplex models follow a turn-by-turn protocol, waiting for explicit end-of-turn signals (e.g., end-pointing/VAD) before responding, which simplifies streaming but adds latency and suppresses natural overlap. They commonly adopt next segment prediction \citep{hara18_interspeech,li-etal-2022-speak,5494991} where the system predicts the agent's response as a segment. Full-duplex models listen and speak simultaneously, modeling overlapping speech, micro-pauses, and background noise to sustain context and deliver timely cues (e.g., continuers, barge-in handling). They either employ the next segment prediction paradigm \citep{arxiv:2503.01174,dgslm,vap} or introduce the next dual token prediction \citep{arxiv:2203.16502, defossez2024moshi} scheme to incorporate the listener's branch.
\paragraph{Faithful Rationales.}
In speech analysis tasks, generating verifiable, evidence-aligned rationale to explain model decisions has been widely studied; such explanations can improve decision reliability and make the reasoning process more transparent. 
Recent work further leverages LLMs, e.g., Brown et al.~\citeyearpar{brown2020languagemodelsfewshotlearners}, together with RL post-training methods such as GRPO \citep{zhang2025grpoleaddifficultyawarereinforcementlearning}, to align and optimize the model's reasoning process and to generate explanations.
This line of work is related to emotion-aware reasoning \citep{wang2026emotionthinkerprosodyawarereinforcementlearning,xu2023secapspeechemotioncaptioning,cheng2024emotionllamamultimodalemotionrecognition,zhang2025classificationspeechemotionreasoning} and mental-state modeling \citep{sclar2024explore}. However, in the full duplex setting, the system needs to make millisecond-level low-latency judgments about speaker behavior, while the computational overhead of explanation generation must be controllable and predictable. To this end, we propose a Graph-of-Thought (GoT) representation inspired by human online perception and incremental reasoning mechanisms, to efficiently model dialogue dynamics and generate high-quality explanations under low-latency constraints.
\paragraph{Conversational Behaviors.}
Interactive behaviors in SDMs have been extensively examined to foster mutual understanding, user engagement, and social connection between human and computer. However, recent studies in SDMs only focus on low-level behaviors (backchannel~\citep{schegloff1982discourse, arvix:2507.23159, arxiv:2503.04721}, turn-taking~\citep{gravano2011turn, duncan1972some, raux2012optimizing}, interruption~\citep{khouzaimi2016reinforcement, marge2022spoken},  continuation~\citep{arvix:2507.23159, arxiv:2503.04721, arora2025talking, nguyen2023generative}, emotion~\citep{arxiv:2508.17623, arvix:2510.07355}, and disfluency~\citep{lian-anumanchipalli-2024-towards-hudm,lian2023unconstrained-udm,ssdm,lian2024ssdm2.0}, which overlook the importance of discourse-level intent that drives these phenomena. To resolve this, we also introduce high-level speech acts (constative, directive, commissive, acknowledgment) ~\citep{jurafsky2025speech} to restore this layer, enabling more interpretable modeling and evaluation for conversational behaviors.

